\documentclass[12pt,letterpaper]{article}
\usepackage{graphicx,textcomp}
\usepackage{natbib}
\usepackage{setspace}
\usepackage{fullpage}
\usepackage{color}
\usepackage[reqno]{amsmath}
\usepackage{amsthm}
\usepackage{fancyvrb}
\usepackage{amssymb,enumerate}
\usepackage[all]{xy}
\usepackage{endnotes}
\usepackage{lscape}
\newtheorem{com}{Comment}
\usepackage{float}
\usepackage{hyperref}
\newtheorem{lem} {Lemma}
\newtheorem{prop}{Proposition}
\newtheorem{thm}{Theorem}
\newtheorem{defn}{Definition}
\newtheorem{cor}{Corollary}
\newtheorem{obs}{Observation}
\usepackage[compact]{titlesec}
\usepackage{dcolumn}
\usepackage{tikz}
\usetikzlibrary{arrows}
\usepackage{multirow}
\usepackage{xcolor}
\newcolumntype{.}{D{.}{.}{-1}}
\newcolumntype{d}[1]{D{.}{.}{#1}}
\definecolor{light-gray}{gray}{0.65}
\usepackage{url}
\usepackage{listings}
\usepackage{color}

\definecolor{codegreen}{rgb}{0,0.6,0}
\definecolor{codegray}{rgb}{0.5,0.5,0.5}
\definecolor{codepurple}{rgb}{0.58,0,0.82}
\definecolor{backcolour}{rgb}{0.95,0.95,0.92}

\lstdefinestyle{mystyle}{
	backgroundcolor=\color{backcolour},   
	commentstyle=\color{codegreen},
	keywordstyle=\color{magenta},
	numberstyle=\tiny\color{codegray},
	stringstyle=\color{codepurple},
	basicstyle=\footnotesize,
	breakatwhitespace=false,         
	breaklines=true,                 
	captionpos=b,                    
	keepspaces=true,                 
	numbers=left,                    
	numbersep=5pt,                  
	showspaces=false,                
	showstringspaces=false,
	showtabs=false,                  
	tabsize=2
}
\lstset{style=mystyle}
\newcommand{\Sref}[1]{Section~\ref{#1}}
\newtheorem{hyp}{Hypothesis}

\title{Problem Set 2}
\date{Due: February 18, 2026}
\author{Applied Stats II}


\begin{document}
	\maketitle
	\section*{Instructions}
	\begin{itemize}
		\item Please show your work! You may lose points by simply writing in the answer. If the problem requires you to execute commands in \texttt{R}, please include the code you used to get your answers. Please also include the \texttt{.R} file that contains your code. If you are not sure if work needs to be shown for a particular problem, please ask.
		\item Your homework should be submitted electronically on GitHub in \texttt{.pdf} form.
		\item This problem set is due before 23:59 on Wednesday February 18, 2026. No late assignments will be accepted.

	\end{itemize}

	\vspace{.25cm}

\noindent We're interested in what types of international environmental agreements or policies people support (\href{https://www.pnas.org/content/110/34/13763}{Bechtel and Scheve 2013)}. So, we asked 8,500 individuals whether they support a given policy, and for each participant, we vary the (1) number of countries that participate in the international agreement and (2) sanctions for not following the agreement. \\

\noindent Load in the data labeled \texttt{climateSupport.RData} on GitHub, which contains an observational study of 8,500 observations.

\begin{itemize}
	\item
	Response variable: 
	\begin{itemize}
		\item \texttt{choice}: 1 if the individual agreed with the policy; 0 if the individual did not support the policy
	\end{itemize}
	\item
	Explanatory variables: 
	\begin{itemize}
		\item
		\texttt{countries}: Number of participating countries [20 of 192; 80 of 192; 160 of 192]
		\item
		\texttt{sanctions}: Sanctions for missing emission reduction targets [None, 5\%, 15\%, and 20\% of the monthly household costs given 2\% GDP growth]
		
	\end{itemize}
	
\end{itemize}

\newpage
\noindent Please answer the following questions:

\begin{enumerate}
	\item
	Remember, we are interested in predicting the likelihood of an individual supporting a policy based on the number of countries participating and the possible sanctions for non-compliance.
	\begin{enumerate}
		\item [] Fit an additive model. Provide the summary output, the global null hypothesis, and $p$-value. Please describe the results and provide a conclusion.

	\end{enumerate}
	
	\item
	If any of the explanatory variables are statistically significant in this model, then:
	\begin{enumerate}
		\item
		For the policy in which nearly all countries participate [160 of 192], how does increasing sanctions from 5\% to 15\% change the odds that an individual will support the policy? (Interpretation of a coefficient)
		\item
		For the policy in which very few countries participate [20 of 192], how does increasing sanctions from 5\% to 15\% change the odds that an individual will support the policy? (Interpretation of a coefficient)
		\item
		What is the estimated probability that an individual will support a policy if there are 80 of 192 countries participating with no sanctions? 
	
	\end{enumerate}
	\item
	Would the answers to 2a and 2b potentially change if we included an interaction term in this model? Why? 
	\begin{itemize}
		\item Perform a test to see if including an interaction is appropriate.
	\end{itemize}
	\end{enumerate}
\newpage
\section*{Problem 1}

\noindent We fit an additive logistic regression model to predict the likelihood of an individual supporting an international environmental policy based on the number of participating countries and the level of sanctions for non-compliance.
\vspace{1em}


\noindent\textbf{R code used:}

\lstinputlisting[language=R, firstline=49, lastline=83]{PS02_JL.R}
\newpage

\noindent\textbf{Model Summary:}
\begin{table}[H] \centering 
	\caption{Logistic Regression: Support for International Environmental Policy} 
	\label{tab:model1} 
	\begin{tabular}{@{\extracolsep{5pt}}lc} 
		\\[-1.8ex]\hline 
		\hline \\[-1.8ex] 
		& \multicolumn{1}{c}{\textit{Dependent variable:}} \\ 
		\cline{2-2} 
		\\[-1.8ex] & choice \\ 
		\hline \\[-1.8ex] 
		countries.L & 0.458$^{***}$ \\ 
		& (0.038) \\ 
		& \\ 
		countries.Q & $-$0.010 \\ 
		& (0.038) \\ 
		& \\ 
		sanctions.L & $-$0.276$^{***}$ \\ 
		& (0.044) \\ 
		& \\ 
		sanctions.Q & $-$0.181$^{***}$ \\ 
		& (0.044) \\ 
		& \\ 
		sanctions.C & 0.150$^{***}$ \\ 
		& (0.044) \\ 
		& \\ 
		Constant & $-$0.006 \\ 
		& (0.022) \\ 
		& \\ 
		\hline \\[-1.8ex] 
		Observations & 8,500 \\ 
		Log Likelihood & $-$5,784.130 \\ 
		Akaike Inf. Crit. & 11,580.260 \\ 
		\hline 
		\hline \\[-1.8ex] 
		\textit{Note:}  & \multicolumn{1}{c}{$^{*}$p$<$0.1; $^{**}$p$<$0.05; $^{***}$p$<$0.01} \\ 
		& \multicolumn{1}{c}{\parbox[t]{7cm}{$^{***}p<0.01$; $^{**}p<0.05$; $^*p<0.1$}} \\ 
	\end{tabular} 
\end{table}

\noindent\textbf{Global Null Hypothesis Test:}

\noindent The likelihood ratio test evaluates the global null hypothesis that all slopes equal zero (i.e., neither countries nor sanctions affect policy support). The test statistic is computed as the difference between the null deviance and residual deviance, which follows a chi-squared distribution with degrees of freedom equal to the number of parameters estimated.

\noindent Results from the global null hypothesis test:
\begin{itemize}
	\item Test Statistic (Chi-squared): $\chi^2 = 215.15$
	\item Degrees of Freedom: $df = 5$
	\item P-value: $p < 2.220446e-16$ (effectively 0)
\end{itemize}

\noindent\textbf{Interpretation and Conclusion:}

\noindent The results strongly reject the global null hypothesis ($\chi^2(5) = 215.15$, $p \approx 0$), indicating that the additive model is statistically significant and the explanatory variables meaningfully improve prediction of policy support.

\noindent Examining the individual coefficients:

\begin{itemize}
	\item \textbf{Countries (Linear term):} The linear contrast for countries is highly significant ($\hat{\beta} = 0.458$, $p < 0.001$). This indicates that as the number of participating countries increases, individuals are significantly more likely to support the policy. The positive coefficient suggests that broader international participation increases policy support.
	
	\item \textbf{Countries (Quadratic term):} The quadratic term for countries is not statistically significant ($\hat{\beta} = -0.010$, $p = 0.794$), suggesting that the relationship between country participation and support is approximately linear rather than curved.
	
	\item \textbf{Sanctions (Linear term):} The linear contrast for sanctions is highly significant and negative ($\hat{\beta} = -0.276$, $p < 0.001$). This indicates that as sanctions increase, individuals are significantly less likely to support the policy. This finding is intuitive: harsher penalties for non-compliance reduce policy support.
	
	\item \textbf{Sanctions (Quadratic term):} The quadratic term for sanctions is also statistically significant and negative ($\hat{\beta} = -0.181$, $p < 0.001$), suggesting a nonlinear relationship where the effect of sanctions accelerates at higher levels.
	
	\item \textbf{Sanctions (Cubic term):} The cubic term for sanctions is statistically significant and positive ($\hat{\beta} = 0.150$, $p < 0.001$), indicating some curvature in the relationship at the highest sanction levels.
\end{itemize}

\noindent\textbf{Conclusion:} The additive logistic regression model is highly significant and explains meaningful variation in policy support. Both the number of participating countries and the level of sanctions significantly affect the probability of supporting international environmental policies. Specifically, individuals are more likely to support policies with broader international participation and less likely to support policies with severe sanctions for non-compliance. The significant polynomial terms for sanctions suggest that the deterrent effect of sanctions is not strictly linear but varies across different sanction levels.


\end{document}
